\documentclass[12pt,a4paper]{report}
\usepackage[]{amsmath,amsthm,amssymb,amscd}
\usepackage[latin1]{inputenc}
%\usepackage{fullpage}

\topmargin 0pt
\advance \topmargin by -\headheight
\advance \topmargin by -\headsep
     
\textheight 246mm
     
\oddsidemargin 0pt
\evensidemargin \oddsidemargin
\marginparwidth 0.5in
     
\textwidth 159mm


%               Theorem Environments
\theoremstyle{definition}
\newtheorem{ex}{}
\newcommand{\pd}[1]{\frac{\partial}{\partial #1}} %\pd{x}
%               Subquestions numbered with letters
\renewcommand{\theenumi}{\textbf{\alph{enumi}}}

%               Maths definitions

\newcommand{\NN}{\ensuremath{\mathbb N}}
\newcommand{\ZZ}{\ensuremath{\mathbb Z}}
\newcommand{\CC}{\ensuremath{\mathbb C}}
\newcommand{\TT}{\ensuremath{\mathbb T}}
\newcommand{\RR}{\ensuremath{\mathbb R}}
\newcommand{\QQ}{\ensuremath{\mathbb Q}}
\newcommand{\g}{\ensuremath{\frak{g}}}
\newcommand{\h}{\ensuremath{\frak{h}}}
\newcommand{\ft}{\ensuremath{\frak{t}}}
\newcommand{\cB}{\mathcal{B}}
\newcommand{\cN}{\mathcal{N}}
\newcommand{\cL}{\mathcal{L}}
\newcommand{\cD}{\mathcal{D}}
%\newcommand{\pd}[1]{\frac{\partial}{\partial #1}} %\pd{x}
%\newcommand{\pd}[1]{\frac{\partial}{\partial #1}} %\pd{x}

\begin{document}
\noindent
{\sffamily\large Marco Zambon 
\hfill     29/03/2011\\Geometr\'ia III, curso 2010-11}
 

%\noindent
%\hrulefill{}\\
\begin{center}\bfseries\large\textsf{Examen Parcial: soluciones    \vspace{-2mm}}
\end{center}

 
%\hrulefill{}\\

 \noindent
\hrulefill{}\\
\textsf{Aqu\'i pongo, para cada problema, una soluci\'on. Est\'a claro que puede haber otras soluci\'ones tambi\'en! \vspace{-2mm}}

\noindent
\hrulefill{}\\

 
\begin{ex}\emph{[8 puntos]    
Considera la variedad topol\'ogica con borde $$M:=\{(x,y):x,y \in \RR,  1 \le \sqrt{x^2+y^2}\le 2\}$$
y considera esta relaci\'on de equivalencia   en $M$: cada punto de $M$ es equivalente s\'olo a s\'i mismo excepto  los puntos del borde $\{(x,y): \sqrt{x^2+y^2}=1\}\cup \{(x,y): \sqrt{x^2+y^2}=2\}$, para los que se cumple
$$(x,y) \sim (2x,2y)$$
para cada $(x,y) \in M$ con $\sqrt{x^2+y^2}=1$.\\
Construye un homeomorfismo del cociente  $M/\sim$ a una superficie en
la lista del teorema de clasificaci\'on de superficies compactas y conexas.\\
}

Construimos un homeomorfismo de $M/\sim$ al toro $S^1 \times S^1$.
Utilizamos la identificaci\'on $\RR^2\cong \CC, (x,y)\mapsto x+i y$ para simplificar\footnote{Para escribir la demonstraci\'on en terminos de $\RR^2$ utiliza $|x+iy|=\sqrt{x^2+y^2}$ para $x,y \in \RR$ y $e^{i\theta}=cos(\theta)+i sin(\theta)$
para $\theta\in \RR$.} 
 la notaci\'on, es decir, tomamos $M=\{z\in \CC:1\le |z|\le 2\}$ y $z\sim 2 z$ para todos $z$ con $|z|=1$.
La aplicaci\'on $$F \colon M \to S^1\times S^1, z\mapsto \left(\frac{z}{|z|},e^{2 \pi i |z|}\right)
$$
es claramente continua. $F$ es sobreyectiva: dados $(w_1,w_2)\in S^1\times S^1$,
tenemos $F(w_1r)=(w_1,w_2)$ donde $r\in [1,2)$ est\'a determinado por $e^{2  \pi i r}=w_2$. Como  $1\le |z|\le 2$ para todos $z\in M$,  tenemos 
\begin{align*}
F(z)=F(z') \Leftrightarrow &
 |z|-|z'|\in \ZZ \text{ y } \frac{z}{|z|}=\frac{z'}{|z'|}\\
\Leftrightarrow &(|z|=|z'| \text{ o bien } |z|=1,|z'|=2 \text{ o bien }|z|=2,|z'|=1) \text{ y }
\frac{z}{|z|}=\frac{z'}{|z'|}\\
\Leftrightarrow &z=z' \text{ o bien } |z|=1, z=\frac{z'}{2} \text{ o bien }|z'|=1, z'=\frac{z}{2}\\
 \Leftrightarrow &
z \sim z'.    
\end{align*}
Ademas $M$ es compacta y el toro $S^1 \times S^1$ es Hausdorff, por lo tanto podemos utilizar el ejercicio 3 de la hoja 1 y concluir que la aplicaci\'on $$\tilde{F} \colon M/\sim \to S^1 \times S^1,$$ determinada por $\tilde{F} \circ \pi = F$ donde $\pi \colon M \to M/\sim$ es la proyecci\'on canonica, es un homeomorfismo.
\end{ex}

\begin{ex}\emph{ [7 puntos]
Sea $P_2(\RR)$ el plano proyectivo y $\TT:=S^1\times S^1$ el toro. > A qu\'e superficie en la lista de superficies compactas y conexas es homeomorfa la suma conexa $$P_2(\RR) \; \sharp \;  \TT\;\;?$$  
Demuestra.\\}

Utilizamos el    teorema de clasificaci\'on de superficies compactas y conexas (sin borde), que clasifica dichas superficies por su orientabilidad y numero de Euler. 
$P_2(\RR) \; \sharp \;  \TT$ no es orientable, porque $P_2(\RR)$ no lo es (contiene una banda de M\"obius). Tenemos $$\chi(P_2(\RR) \; \sharp \;  \TT)=
\chi(P_2(\RR))+ \chi(\TT)-2=1+0-2=-1.$$
Por lo tanto $P_2(\RR) \; \sharp \;  \TT$ es homeomorfo 
%a una suma conexa de varias copias de $P_2(\RR)$ con $\chi=-1$, es decir, a
 $P_2(\RR)  \; \sharp \;  P_2(\RR)  \; \sharp \;  P_2(\RR)$, ya que esta no es orientable y tiene numero de Euler $-1$.
\end{ex}
 

\begin{ex} \emph{ [8 puntos]
> Es verdad que $\RR$, con su topolog\'ia canonica, admite un numero infinito de estructuras  de variedad diferencial (distintas entre s\'i)? Demu\'estralo.\\
}

S\'i, es verdad. Para cada $n\in \NN_{\ge 0}$ sea ${\cal A}_{2n+1}$
el atlas diferencial maximal que contiene el homeomorfismo $\phi_{2n+1} \colon \RR \to \RR, x \mapsto x^{2n+1}$. Los ${\cal A}_{2n+1}$ son distintos dos a dos: dados $n,m \in \NN_{\ge 0}$, digamos con $n<m$, tenemos que $\phi_{2n+1} \notin {\cal A}_{2m+1}$ porque el cambio de coordenadas $\phi_{2n+1}\circ(\phi^{2m+1})^{-1} \colon \RR \to \RR$ es $x^{\frac{2n+1}{2m+1}}$ y por lo tanto no es diferenciable en zero. 
\end{ex}





\begin{ex}[7 puntos]\emph{
Sea $M$ una variedad diferencial de dimensi\'on $n$, $S$  una subvariedad de $M$ de dimensi\'on $k$, y $p$ un punto de $S$. Demuestra que existe un entorno abierto $U$ de $p$ en $M$ y una aplicaci\'on diferenciable $F \colon U \to \RR^{n-k}$ tal que
$$U\cap S= F^{-1}(0).$$\\
}

Por definici\'on de subvariedad, existe una carta $(U,\phi)$ de la variedad diferencial $M$ adaptada a $S$ y con $p\in U$. Es decir, $U$ es un entorno abierto de $p$ en $M$ y $\phi \colon U \to \phi(U)$ es un difeomorfismo a un abierto de $\RR^n$   tal que $\phi(U\cap S)=\phi(U)\cap (\RR^{k}\times \{0\})$, donde $\{0\}$ es la origen en $\RR^{n-k}$. Considera la aplicaci\'on $$F:= \pi \circ \phi \colon U \to \RR^{n-k}$$ donde $\pi \colon \RR^{n}\to \RR^{n-k}$ es la proyecci\'on a las \'ultimas $n-k$ componentes. $F$ es diferenciable porque bajo la carta $\phi$ corresponde  a $\pi$, que es una aplicaci\'on diferenciable entre abiertos del espacio euclideo, o alternativamente porque $F$ es la composici\'on de dos aplicaci\'ones diferenciables.
Ademas \begin{align*}
F^{-1}(0)&=\phi^{-1}(\pi^{-1}(0))\\
&=\phi^{-1}(\RR^{k}\times \{0\})\\
&=\phi^{-1}(\phi(U)\cap (\RR^{k}\times \{0\})    )\\
&= U\cap S.   
\end{align*}
\end{ex}

 

%\begin{ex}  
%Demuestra que 
%$$T \;\;\sharp \;\; C\cong T-(D_1 \cup D_2)$$
%donde $T=S^1\times S^1$ es el toro, $C=S^1\times [0,1]$ el cilindro, y $D_i\subset T$ (por $i=1,2$) son discos cerrados tal que $D_1\cap D_2=\emptyset$.
%\end{ex} 






 \end{document}
 %%%%%%%
 \begin{ex}

\end{ex}

 \begin{ex}
\begin{enumerate}
\item 
\end{enumerate}
\end{ex}
 
 
 
 